\documentclass[10pt,a4paper]{amsart}
\usepackage[margin=19mm]{geometry}
\usepackage{amsmath,amssymb,mathtools}
\usepackage[scr=rsfs,cal=euler]{mathalfa}
\usepackage{xfrac}
\usepackage[unicode,hidelinks,bookmarks=false]{hyperref}
\newcommand{\ie}{i.e.\ }
\newcommand{\cf}{cf.\ }
\newcommand{\resp}{respectively\ }
\newcommand{\Subsection}{\subsection}
\providecommand{\nobreakdash}{\nobreak\hbox{-}}
\setlength{\emergencystretch}{2em}
\multlinegap0pt
\usepackage{braket}
\usepackage{amssymb}
\usepackage{mathtools}
\usepackage{float}
\usepackage{caption}
\newtheorem{thm}{Theorem}
\newtheorem{claim}{Claim}
\title{On graphs with large third eigenvalue}
\author{Giacomo Leonida, Sida Li}
\date{Source: arXiv:2501.02563v4 (CC BY 4.0)}
\begin{document}
\maketitle

\noindent\emph{Source and licence.} On graphs with large third eigenvalue. Version arXiv:2501.02563v4 (CC BY 4.0), distributed by arXiv under the Creative Commons Attribution 4.0 International licence, http://creativecommons.org/licenses/by/4.0/, as recorded in the arXiv licence metadata for this exact version and shown on the abstract page https://arxiv.org/abs/2501.02563v4. Full licence notice: \texttt{licenses/arxiv-2501.02563-CC-BY-4.0.txt}.\par\smallskip

\noindent\textit{Editorial scope.} Counted unit: the article's thm lambda\_n<lambda\_n-1 (source0.tex lines 445--451). The article's complete original proof of it is delivered with it (source0.tex lines 452--499, one \texttt{proof} environment of 48 lines). No mathematics is written, repaired or re-derived: the statement and the proof are the article's own lines, in the article's own notation, and no retained line is substituted or re-flowed.  No further source-local context is needed: every cross-reference of every retained line resolves inside the two delivered spans.  Nothing was omitted from any retained span, so the package carries no omission record.  No retained line contains a citation, so the wrapper carries no bibliography and no bibliography entry.  No retained line contains a figure environment, an image-inclusion command or drawing code, so no image is delivered and the package declares no asset dependency.  Editorial formatting only, in addition: an AMS \texttt{amsart} preamble that restates the article's own declarations and macros used by the retained lines, namely the environments \texttt{thm}, \texttt{claim} on separate counters, together with the standard packages those lines need.  The retained environments print in this document as Theorem 1, Claim 1 in order of appearance, whereas the article prints the same items with the numbering of its own full document.  The complete native source of the article as distributed in the arXiv e-print for this exact version is bundled unchanged as \texttt{sources/171176/source0.tex}.
\begin{thm} \label{lambda_n<lambda_n-1} Suppose $M\in \mathcal{S}_n$ is a simple minimiser. Then, up to reordering rows and columns:
\begin{equation*}
    M = \begin{pmatrix} O_{P} & J_{P,N} & X \\ J_{N,P} & O_{N} & Y \\ X^T & Y^T & W
        \end{pmatrix}
\end{equation*}
where $O_m$ denotes the $m\times m$ matrix of zeros, $J_{m,l}$ denotes the $m\times l$ matrix of ones and $X,Y,W$ are $P\times Z$, $N\times Z$ and $Z\times Z$ matrices with $[0,1]$ entries respectively, with $\frac{1}{\sqrt{P}}X^T\bold{j}_P = \frac{1}{\sqrt{N}}Y^T\bold{j}_N$. If $\bold{w}$ denotes the unit eigenvector of $M$ corresponding to $\lambda_{n-1}$, then $\bold{w}^T = \begin{pmatrix} \frac{\bold{j}_P^T}{\sqrt{2P}} & -\frac{\bold{j}_N^T}{\sqrt{2N}} & \bold{0}_{Z}^T \end{pmatrix}$ and $\lambda_{n-1} = -\sqrt{PN}$.
\end{thm}

\begin{proof}
    Suppose that $M$ is a simple minimiser with eigenvalues $\lambda_n(M)<\lambda_{n-1}(M)$ and corresponding orthonormal eigenvectors $\bold{v,w}$. Suppose that $E$ is an $n\times n$ matrix such that $\exists \eta>0$ with $M+\varepsilon E \in \mathcal{S}_n$ for all $0<\varepsilon<\eta$. Then, by the min-max principle:
    \begin{alignat*}{1}
        \lambda_{n-1}(M+\varepsilon E) - \lambda_{n-1}(M) & \le \max_{\substack{\bold{x} \in \braket{\bold{v,w}} \\ ||\bold{x}|| = 1}} \left(\bold{x}^{T}(M+\varepsilon E)\bold{x} \right) - \lambda_{n-1}(M) \\
        & = \max_{\substack{\alpha, \beta \in [-1,1]\\ \alpha^2 + \beta^2 = 1}} \left(\alpha^2 (\lambda_{n}(M)-\lambda_{n-1}(M)) + \varepsilon(\alpha\bold{v} + \beta \bold{w})^{T}  E (\alpha\bold{v} + \beta \bold{w})\right)\\ & = \max_{\alpha\in [-1,1]} \left(\alpha^2 (\lambda_n(M)-\lambda_{n-1}(M) + \varepsilon(S-R)) + 2\varepsilon \alpha\sqrt{1-\alpha^2} Q + \varepsilon R\right),
    \end{alignat*}
    where $S = \bold{v}^{T}E\bold{v}$, $Q = \bold{v}^{T}E\bold{w}$ and $R = \bold{w}^{T}E\bold{w}$. Provided that $\lambda_{n-1}(M)-\lambda_{n}(M)>\varepsilon (S-R)$ (which holds for $\varepsilon$ sufficiently small), we have:
    \begin{alignat*}{1}
        \lambda_{n-1}(M+\varepsilon E) - \lambda_{n-1}(M) & = \max_{\alpha\in [-1,1]} \left(\alpha^2 (\lambda_n(M)-\lambda_{n-1}(M) + \varepsilon(S-R)) + 2\varepsilon \alpha\sqrt{1-\alpha^2} Q + \varepsilon R\right)\\ & \le  \max_{\alpha\in [-1,1]} \left(\alpha^2 (\lambda_n(M)-\lambda_{n-1}(M) + \varepsilon(S-R)) + 2\varepsilon \alpha Q + \varepsilon R\right) \\ & = \frac{\varepsilon R(\lambda_{n}(M)-\lambda_{n-1}(M) + \varepsilon(S-R)) - \varepsilon^2 Q^2}{\lambda_n(M)-\lambda_{n-1}(M) + \varepsilon(S-R)}. \tag{*}
    \end{alignat*}
    Here, we chose $\alpha$ which maximises the above quadratic.
    \begin{claim} If $E$ is a matrix such that $M+\varepsilon E \in \mathcal{S}_n$ for all $\varepsilon$ sufficiently small, then $R \ge 0$. 
    \end{claim}
    \begin{proof}
    If by way of contradiction $R < 0$ then (*) is negative for all $\varepsilon > 0$ sufficiently small. Thus, for $\varepsilon$ sufficiently small we have $\lambda_{n-1}(M+\varepsilon E)<\lambda_{n-1}(M)$, contradicting the minimality of $\lambda_{n-1}(M)$ amongst elements of ${\mathcal{S}}_n$. \end{proof}
    Let $W_{+} = \{i:w_i >0\}$, $W_{0} = \{i:w_i =0\}$ and $W_{-} = \{i:w_i<0\}$ and let $P = |W_{+}|$, $N = |W_{-}|$ and $Z = |W_{0}|$. WLOG we may reorder the vertices of the graph (i.e. rows and columns of $M$) such that $W_{+} = \{1,\dots, P\}$, $W_{-} = \{P+1,\dots, P+N\}$ and $W_{0} = \{P+N+1,\dots, P+N+Z\}$. 
    
    If by way of contradiction there exists $i,j$ with $(i,j) \in W_{+}\times W_{+} \cup W_{-}\times W_{-}$ and $M_{ij} > 0$, then consider $E$ of the form:
    \begin{equation*}
        E_{xy} = \begin{cases} -1 & \{x,y\} = \{i,j\} \\ 0 & \text{otherwise.}
        \end{cases}
    \end{equation*}
    For such $E$, $R = \bold{w}^{T}E\bold{w} = -2w_iw_j < 0$, so by the claim above, $M$ is not minimal (yielding a contradiction). Hence, if $i\ne j$ are such that $(i,j) \in W_{+}\times W_{+} \cup W_{-}\times W_{-}$, then $M_{ij} = 0$.  
    
    Similarly, if $i,j$ are such that $(i,j) \in W_{+}\times W_{-} \cup W_{-}\times W_{+}$ and $M_{ij}<1$ then considering:
    \begin{equation*}
        E_{xy} = \begin{cases} 1 & \{x,y\} = \{i,j\} \\ 0 & \text{otherwise,}
        \end{cases}
    \end{equation*}
    we again have $R = \bold{w}^{T}E\bold{w} = 2w_iw_j < 0$ giving a contradiction. Thus, if $i,j$ are such that $(i,j) \in W_{+}\times W_{-} \cup W_{-}\times W_{+}$, then $M_{ij}=1$. Hence, $M$ must have the form:
    \begin{equation*}
        M = \begin{pmatrix} O_{P} & J_{P,N} & X \\ J_{N,P} & O_{N} & Y \\ X^T & Y^T & W
        \end{pmatrix},
    \end{equation*}
    where $O_m$ denotes the $m\times m$ matrix of zeros, $J_{m,l}$ denotes the $m\times l$ matrix of ones and $X,Y,W$ are $P\times Z$, $N\times Z$ and $Z\times Z$ matrices with $[0,1]$ entries respectively. Writing $\bold{w}^T = \begin{pmatrix}\bold{w}^T_P & \bold{w}^T_N & \bold{0}_{Z}^T\end{pmatrix}$, as $\bold{w}$ is the penultimate eigenvector of $M$, we require:
    \begin{align*}
        J_{P,N}\bold{w}_{N} & = \lambda_{n-1} \bold{w}_{P} \\
        J_{N,P}\bold{w}_{P} & = \lambda_{n-1} \bold{w}_{N} \\
        X^T\bold{w}_{P} + Y^T\bold{w}_{N} & = \bold{0}_{Z}.
    \end{align*}
    Namely, the first two equations give $\bold{w}_P = x\bold{j}_P$ and $\bold{w}_N = -y\bold{j}_N$ for some $x,y>0$. Solving for $x,y$ and $\lambda_{n-1}$ gives:
    \begin{equation*}
        \lambda_{n-1} = -\sqrt{PN},\quad \bold{w}^T = \begin{pmatrix}
            \frac{\bold{j}_P^T}{\sqrt{2P}} & -\frac{\bold{j}_N^T}{\sqrt{2N}} & \bold{0}_{Z}^T
        \end{pmatrix},
    \end{equation*}
    and $\frac{X^{T}\bold{j}_{P}}{\sqrt{P}} = \frac{Y^{T}\bold{j}_{N}}{\sqrt{N}}$. 
\end{proof}

\end{document}
